\section{The problem as optimization}\label{sec:problem}

\inbrief{Each pair of hard blocks must sit left of, right of, above or below one
another, so with $m$ pairs the legal placements form a union of up to $4^m$ convex
pieces, and a search that moves one block at a time can leave its piece only where
two pieces touch. The score averages the most crowded grid cells and assigns pins to
cells, so it has kinks and jumps where gradients fail, and one exact evaluation takes
seconds.}

\subsection{Variables and constraints}\label{sec:geometry}\label{sec:geometry-analysis}

The canvas is the rectangle $[0,L_x]\times[0,L_y]$. Object $i$, a hard macro or a
soft cluster, has center $p_i=(x_i,y_i)$, width $w_i$ and height $h_i$; sizes and
orientations are fixed. A placement $p$ stacks all centers into one vector, of
2,156--5,536 coordinates on the IBM designs. Every object must lie inside the
canvas,
\begin{equation}
  \tfrac{w_i}{2}\le x_i\le L_x-\tfrac{w_i}{2},\qquad
  \tfrac{h_i}{2}\le y_i\le L_y-\tfrac{h_i}{2},
  \label{eq:containment}
\end{equation}
and fixed objects keep their positions (no IBM design has any). Soft clusters may
overlap anything, but two hard macros $i\neq j$ must be separated horizontally or
vertically: at least one of four linear inequalities must hold,
\begin{equation}
\Bigl(x_i+\tfrac{w_i}{2}\le x_j-\tfrac{w_j}{2}\Bigr)\lor
\Bigl(x_j+\tfrac{w_j}{2}\le x_i-\tfrac{w_i}{2}\Bigr)\lor
\Bigl(y_i+\tfrac{h_i}{2}\le y_j-\tfrac{h_j}{2}\Bigr)\lor
\Bigl(y_j+\tfrac{h_j}{2}\le y_i-\tfrac{h_i}{2}\Bigr).
\label{eq:disjunction}
\end{equation}
Touching edges are legal. Choosing one of the four options for every pair turns this
either-or rule into ordinary linear inequalities, which gives the structure below.
Let $\mathcal B$ be the set of placements that satisfy
\eqref{eq:containment} and the fixed positions, $\mathcal E$ the set of unordered
hard-macro pairs with $m=|\mathcal E|$, and $H_{e,b}$ the closed halfspace given by
branch $b\in\{1,2,3,4\}$ of \eqref{eq:disjunction} for pair $e$.

\begin{proposition}[Finite polyhedral decomposition]\label{prop:geometry}
The feasible set is a union of at most $4^m$ compact convex polyhedra:
\begin{equation}
 \mathcal F
 =\mathcal B\cap\bigcap_{e\in\mathcal E}\bigcup_{b=1}^4 H_{e,b}
 =\bigcup_{\sigma\in\{1,2,3,4\}^{\mathcal E}} P_\sigma,
 \qquad
 P_\sigma=\mathcal B\cap\bigcap_{e\in\mathcal E}H_{e,\sigma(e)}.
 \label{eq:polyhedral-union}
\end{equation}
The polyhedra may overlap or be empty.
\end{proposition}
\noindent Proof in Appendix~\ref{app:proofs}. Each piece $P_\sigma$ fixes, for every
pair, which side of the other each macro lies on (\cref{fig:feasible}); IBM01's 246
hard macros form $m=30{,}135$ pairs, so up to $4^{30{,}135}$ pieces. Inside a piece the constraints are linear and
convex, so continuous optimization works, but changing a pair's arrangement means
moving to another piece, which a path that never overlaps can do only where two
pieces touch. Along one coordinate, with the others fixed, the legal values form a
finite union of closed intervals (Appendix~\ref{app:proofs}): a one-coordinate
search sees gaps where other blocks stand in the way.

\subsection{The score}\label{sec:score}\label{sec:tail-analysis}

The competition proxy, introduced with Circuit Training \citep{mirhoseini2021graph}
and computed by the TILOS MacroPlacement evaluator \citep{macroplacement}, is
\begin{equation}
  \proxy(p)=W(p)+\tfrac12 D(p)+\tfrac12 C(p).
  \label{eq:canonical}
\end{equation}
Appendix~\ref{app:objective} gives every constant.

\paragraph{Wirelength.} The half-perimeter wirelength (HPWL) of a net $n$ is half
the perimeter of the smallest rectangle around its pins, a standard estimate of the
wire needed to connect them. With each pin $q$ at a fixed offset from its object's
center and net weights $\omega_n$,
\[
 \mathrm{HPWL}_n=\bigl(\max_{q\in n}x_q-\min_{q\in n}x_q\bigr)
 +\bigl(\max_{q\in n}y_q-\min_{q\in n}y_q\bigr),
 \qquad
 W=\frac{\sum_n\omega_n\,\mathrm{HPWL}_n}{(L_x+L_y)\sum_n\omega_n}.
\]

\paragraph{Tail averages.} Density and congestion use a grid of $G$ cells (810 to
2,255 per design). For $N$ cell values $z\in\mathbb R^N$, let $A_k(z)$ be the mean of
the $k$ largest: \emph{the average of the worst $k$ cells}. With $\rho_b$ the total
object area in cell $b$ divided by the cell area (it exceeds 1 where objects
overlap), $D=\tfrac12A_{\lfloor0.1G\rfloor}(\rho)$ with $N=G$, so the raw top-decile
mean enters $\proxy$ with weight $\tfrac14$. \emph{Routing congestion} compares
wiring demand with the routing tracks available: the evaluator maps each pin to the
grid cell that contains it, routes each net between those cells with simple L- and
T-shaped paths, divides each cell's horizontal and vertical demand by its capacity,
spreads it evenly over the cell and up to two neighbors on each side along one axis,
and adds a fixed allowance for macros covering the cell. With $u\in\mathbb R^{2G}$
the horizontal and vertical values, $C=A_{\lfloor0.05\cdot2G\rfloor}(u)$ with
$N=2G$, the mean of the worst 5\%.

\begin{lemma}[Tail representation and stability]\label{lem:tail}
For $z\in\mathbb R^N$, an integer $1\le k\le N$ and $(a)_+=\max\{a,0\}$,
\begin{equation}
 A_k(z)
 =\max_{\substack{0\le\beta_j\le1/k\\\sum_j\beta_j=1}}\beta^\top z
 =\min_{\theta\in\mathbb R}\Bigl\{\theta+
 \frac1k\sum_{j=1}^N(z_j-\theta)_+\Bigr\}.
 \label{eq:tail-representation}
\end{equation}
$A_k$ is convex and piecewise affine, and for any $\delta\in\mathbb R^N$,
\begin{equation}
 |A_k(z+\delta)-A_k(z)|
 \le\min\Bigl\{\|\delta\|_\infty,\
          \frac{\|\delta\|_2}{\sqrt{k}},\
          \frac{\|\delta\|_1}{k}\Bigr\}.
 \label{eq:tail-stability}
\end{equation}
\end{lemma}
\noindent The threshold form is the conditional value-at-risk formula
\citep{rockafellar2000cvar}; proof in Appendix~\ref{app:proofs}. For the solver, the
tails are convex in the cell values but kinked wherever the set of worst cells
changes, and they are not convex in block positions, because cell values depend on
positions nonlinearly. A move that changes the cells by $\delta$ moves a tail by at
most $\|\delta\|_1/k$, so a move that touches few cells shifts the crowding terms
little: the first sampled IBM01 move changed 12 of 1,845 cells and the density
component $D=\tfrac12A_k$ by 0.0025, within its bound $\tfrac12\|\delta\|_1/k=0.0153$
(\measurement). The same locality is what \cref{sec:cheap} exploits.

\subsection{Why it is hard}\label{sec:hard}

\begin{itemize}[leftmargin=1.4em,topsep=0.2em,itemsep=0.2em]
  \item \textbf{Nonconvex and disconnected:} good placements in different pieces can
        be separated by overlapping states that no legal path crosses.
  \item \textbf{Piecewise and discontinuous:} pin extremes, clamps and top-$k$
        selection leave kinks, and because the congestion estimate assigns pins to
        grid cells, the score jumps whenever a pin crosses a grid line; its routing
        part has zero gradient almost everywhere.
  \item \textbf{Expensive:} a full evaluation with the reference evaluator took
        1.4\,s on IBM01 and 27\,s on IBM10 (\measurement), so an hour buys about
        2,600 and 130 of them.
  \item \textbf{Large and conflicting:} thousands of coordinates interact through
        7,269--83,132 nets, and packing connected blocks shortens wires but narrows
        routing channels (\cref{fig:placement-tradeoff}); weighted congestion is
        about two-thirds of the score in every development record (\historical).
\end{itemize}

\begin{figure}[tbp]
\centering
\begin{subfigure}[b]{0.6\linewidth}
\centering
\begin{tikzpicture}[x=0.73cm,y=0.53cm,font=\small,
  block/.style={draw=blue!65,fill=blue!8,line width=0.6pt},
  route/.style={draw=orange,line width=0.7pt},
  connection/.style={draw=blue,line width=1.3pt}]
  \begin{scope}
    \node[font=\small\bfseries] at (3,6.55) {Compact};
    \draw[rulegray] (0,0) rectangle (6,6);
    \path[fill=orange!10] (0,2.65) rectangle (6,3.35);
    \draw[block] (0.9,3.35) rectangle (5.1,4.85);
    \draw[block] (0.9,1.15) rectangle (5.1,2.65);
    \node at (3,4.1) {Macro A};
    \node at (3,1.9) {Macro B};
    \foreach \y in {2.8,3.0,3.2} {\draw[route] (0.15,\y) -- (5.85,\y);}
    \draw[connection] (3,2.65) -- (3,3.35);
    \fill[blue] (3,2.65) circle (1.5pt);
    \fill[blue] (3,3.35) circle (1.5pt);
    \node[align=center,font=\footnotesize] at (3,-0.72)
      {Shorter A--B connection;\\narrower channel};
  \end{scope}
  \begin{scope}[xshift=5.2cm]
    \node[font=\small\bfseries] at (3,6.55) {Spread out};
    \draw[rulegray] (0,0) rectangle (6,6);
    \path[fill=orange!10] (0,2.0) rectangle (6,4.0);
    \draw[block] (0.9,4.0) rectangle (5.1,5.5);
    \draw[block] (0.9,0.5) rectangle (5.1,2.0);
    \node at (3,4.75) {Macro A};
    \node at (3,1.25) {Macro B};
    \foreach \y in {2.4,3.0,3.6} {\draw[route] (0.15,\y) -- (5.85,\y);}
    \draw[connection] (3,2.0) -- (3,4.0);
    \fill[blue] (3,2.0) circle (1.5pt);
    \fill[blue] (3,4.0) circle (1.5pt);
    \node[align=center,font=\footnotesize] at (3,-0.72)
      {Longer A--B connection;\\wider channel};
  \end{scope}
  \draw[connection] (0,-1.95) -- (0.7,-1.95);
  \node[anchor=west,font=\footnotesize] at (0.8,-1.95) {A--B connection};
  \draw[route] (5.2,-1.95) -- (5.9,-1.95);
  \node[anchor=west,font=\footnotesize] at (6.0,-1.95) {other nets};
\end{tikzpicture}
\caption{Why the terms compete (schematic).}
\label{fig:placement-tradeoff}
\end{subfigure}\hfill
\begin{subfigure}[b]{0.38\linewidth}
\centering
\begin{tikzpicture}[x=0.86cm,y=0.92cm,font=\footnotesize]
  \fill[blue!14] (0,0) rectangle (1.8,4.2);
  \fill[blue!14] (4.2,0) rectangle (6,4.2);
  \fill[orange!22,opacity=0.75] (0,3.0) rectangle (6,4.2);
  \fill[orange!22,opacity=0.75] (0,0) rectangle (6,1.2);
  \draw[rulegray] (0,0) rectangle (6,4.2);
  \draw[dashed,ink,line width=0.6pt,fill=white] (1.8,1.2) rectangle (4.2,3.0);
  \draw[ink,fill=rulegray!80,line width=0.6pt] (2.3,1.6) rectangle (3.7,2.6);
  \node at (3.0,2.1) {$j$};
  \draw[blue!70!black,fill=blue!30,line width=0.6pt] (0.4,3.2) rectangle (1.4,4.0);
  \node at (0.9,3.6) {$i$};
  \node at (0.9,2.1) {left};
  \node at (5.1,2.1) {right};
  \node at (3.0,3.6) {above};
  \node at (3.0,0.6) {below};
  \draw[<->,muted] (1.8,-0.22) -- node[below]{$w_i+w_j$} (4.2,-0.22);
  \draw[<->,muted] (6.22,1.2) -- node[below,rotate=90]{$h_i+h_j$} (6.22,3.0);
\end{tikzpicture}
\caption{Where $i$'s center may go ($j$ fixed).}
\label{fig:feasible}
\end{subfigure}
\caption{(a)~More separation leaves room for wires crossing between two blocks but
lengthens their connection (schematic; macro interiors obstruct routing). (b)~Each
branch of \eqref{eq:disjunction} confines $i$'s center to a halfplane left of, right
of, below or above $j$; their union, the plane minus the dashed box, is not convex,
and corners satisfy two branches ($i$ is drawn in the upper-left corner).}
\label{fig:problem-pictures}
\end{figure}
